\stagechapter{L}{LLM feasibility}

\noindent\textbf{What this stage covers.} Whether and how a language model can take part in the coding: what task it is given, how the prompt and parameters are fixed, how its output is validated against human labels, and what may be claimed. The draft thesis plan uses the model as a screener with prompts that favour recall: staged prompts (relevant spans, then stance for detected harms, then protection details), each unit run three times, two paraphrased prompts to check that headline counts hold, a preference for an open-weight model if it performs close to a hosted one, raw output kept apart from adjudicated labels, and manual review of all retained harm mentions and potential mismatches. Acceptance thresholds and the model's role are decision D9 (open).

\section{Lineage: from trained classifiers to prompted models}
\pa{wilson2016} created OPP-115 and trained baseline classifiers; \pa{harkous2018} built Polisis on it, evaluated by holding out policies and comparing machine labels with aggregated human labels (accuracy and $\kappa$ per task), and found weakness on retention and user choice, the categories closest to consent quality. \pa{rodriguez2024} replaces the trained pipeline with prompts to GPT-4 Turbo and Llama 2 on the MAPP and OPP-115 corpora: tuning on 33 policies, a control set of 31, cumulative prompt changes (adding the annotation guidelines' definitions was the largest single gain, F1 .848 to .887 \pc{rodriguez2024}{4}{the prompt was augmented to include the internal definitions used for manual annotations in the MAPP dataset}), temperature 0 and a fixed seed, final F1 .935 on the control set \pc{rodriguez2024}{7}{This configuration yields an accuracy of 0.916, a recall of 0.976, a precision of 0.898, and an F1 score of 0.935 on the control set of the MAPP corpus} but Llama 2 weaker (F1 .882 on MAPP, .753 on OPP-115). It also measured determinism: 89.5\,\% identical outputs at temperature 0 \pc{rodriguez2024}{6}{we observed 52 discrepancies (i.e., different responses compared to the typical answer), which means 89.5\% consistency in ChatGPT responses} against 59.6\,\% at default temperature, and found that a persona instruction and over-fine task decomposition can lower accuracy. Its evaluation unit is coarse (a data type present anywhere in the policy) and it reports no chance-corrected agreement.

\section{Workflows and best-practice checklists}
\begin{itemize}
\item \pa{pangakis2023}: a five-step, human-first workflow (codebook; experts label a random subset; LLM labels a tuning subset with the same codebook; refine; test once on the untouched held-out labels), because the LLM prompt \emph{is} the shared codebook \pc{pangakis2023}{6}{Critically, both the human coders and the LLM should use the same codebook for annotation, where the codebook serves as the LLM’s prompt instructions}. Across 27 tasks the median F1 was .707 (range .059--.969), and for about a third of the tasks precision or recall fell below .5, so validation must be task by task. A \emph{consistency score} (share of repeated runs at temperature above 0 that match the modal label) flagged edge cases: accuracy was 19.4 points higher for fully consistent texts.
\item \pa{tornberg2024}: nine points. Fixed and documented model version \pc{tornberg2024}{3}{use a fixed version of the LLM throughout the project, document the version, and ensure that the model will be available for future use}, systematic coding procedure with the prompt developed in the same loop as the guidelines, a \emph{prompt codebook} (definitions, examples, ambiguous-case rules plus the exact prompt, model and parameters), structured prompts with an explicit ``uncertain'' option, validation \emph{after} the prompt is final \pc{tornberg2024}{10}{it is not acceptable to use the validation data to improve the prompts, as this may lead to falsely reporting higher precision} with at least 20--30 examples per category as a practical minimum, explicit parameters even when default, ethics and law (no confidential data to an API without care), prompt-stability checks with paraphrases \pc{tornberg2024}{13}{create several paraphrases of the prompt and run the analysis for a subset of the data} and Krippendorff's $\alpha$ among outputs, and awareness of training-data contamination.
\item \pa{baumann2025}: ``LLM hacking'' \pc{baumann2025}{1}{We call this phenomenon where configuration choices lead to incorrect conclusions LLM hacking}, incorrect conclusions caused by researcher choices of model, prompt, decoding parameters and output mapping. Across 37 replicated tasks and 18 models, a false positive could be produced for 94.4\,\% of null hypotheses by choosing among configurations, and the risk of an incorrect conclusion in a single plausible configuration was about 31\,\% for the best 70B models and up to 50\,\% for 1B models; findings near a significance threshold are most fragile; 100 human labels controlled false positives better than 100{,}000 LLM labels. Recommendations: pre-register all configuration choices \pc{baumann2025}{17}{Pre-registration of all LLM configuration choices, combined with reporting results from all pre-registered configurations, thus makes intentional LLM hacking more difficult}, report every model, prompt and parameter tried, collect as many expert labels as possible, and do not use human agreement alone to decide whether to automate \pc{baumann2025}{34}{Do not rely on human annotator agreement to decide whether annotations should be automated or not}.
\item \pa{gilardi2023} popularised zero-shot annotation (four datasets, 6{,}183 texts, accuracy above crowd-workers). Its ``intercoder agreement'' for ChatGPT is agreement between two runs of the same model \pc{gilardi2023}{6}{intercoder agreement refers to the percentage of instances for which both annotators in a given group report the same class}, not independent coders, and its gold standard keeps only texts on which both trained annotators agreed; later work (Pangakis, Törnberg) treats it as the baseline to correct.
\item \pa{mayring2025}: a field test of ChatGPT as a ``qualitative content analysis according to Mayring'' (naive and precise prompts, a re-run, six interviews) found incorrect procedures, categories not independent, wrong frequencies and a re-run in which ``the coding agreement test must be considered completely failed''; it also warns that sending interviews to ChatGPT needs explicit written consent. It concerns inductive category formation, not deductive coding with a closed codebook.
\item \pa{aoir2019} (companion 6.1) asks that the model version, retraining choices, false positives and negatives and the human override be documented, and warns about taking automated inferences as base truth.
\end{itemize}

\begin{table}[t]
\centering\small
\caption{LLM-annotation sources: task, validation and what each adds.}
\label{tab:L-sources}
\begin{tabularx}{\textwidth}{@{}P{2.6cm}P{3.1cm}P{3.3cm}Y@{}}
\toprule
\textbf{Source} & \textbf{Task / data} & \textbf{Validation} & \textbf{Takeaway}\\
\midrule
\pa{harkous2018} & OPP-115 categories and attributes & accuracy, $\kappa$ vs aggregated human labels; 50 held-out policies & chance-corrected metric per task; lenient reference hides weak classes\\
\pa{rodriguez2024} & data-type extraction, MAPP and OPP-115 & tuning set vs control set; F1; determinism runs & log each prompt change; definitions in prompt; temperature 0 and seed\\
\pa{pangakis2023} & 27 social-science tasks & held-out expert labels; consistency score & human-first workflow; task-by-task validation\\
\pa{tornberg2024} & best-practice essay & validation after prompt freeze; $\alpha$ among paraphrases & checklist; prompt codebook\\
\pa{baumann2025} & 37 tasks, 18 LLMs & human labels vs configurations & pre-register; report all configurations\\
\pa{gilardi2023} & four datasets, three annotator groups & accuracy; run-to-run agreement & why self-consistency is not reliability\\
\pa{mayring2025} & inductive qualitative content analysis & re-run, sample solution & do not use for inductive coding\\
\bottomrule
\end{tabularx}
\end{table}

\begin{notebox}[Points of caution found in the cards]
(1) Agreement of a model with itself measures consistency, not validity or independent reliability (\pa{gilardi2023}, \pa{mayring2025}). (2) F1 against one reference is not comparable with a human--human $\alpha$ (\pa{rodriguez2024}, \pa{pangakis2023} report no chance-corrected statistic). (3) Gold standards that drop disagreements overstate performance (\pa{gilardi2023}, \pa{baumann2025}). (4) Results depend on model snapshots that disappear; prefer a self-hosted open-weight model when quality allows (\pa{tornberg2024}). (5) Near-threshold findings are fragile (\pa{baumann2025}).
\end{notebox}

\begin{protocolbox}[Stage-L checklist assembled from the corpus]
\begin{itemize}[leftmargin=1.3em]
\item Write the decision first: model and version, prompt, temperature, top-$p$, seed, output mapping, selection rule, thresholds, \emph{before} looking at validation labels (\pa{baumann2025}, \pa{tornberg2024}).
\item Split human-labelled material into a tuning set and an untouched validation set; pilot apps never enter validation (as the draft says).
\item Put the codebook definitions into the prompt and archive prompt and codebook together (\pa{pangakis2023}, \pa{rodriguez2024}).
\item Measure determinism (repeat runs), prompt stability (paraphrases, $\alpha$) and a consistency score; keep all raw outputs (\pa{rodriguez2024}, \pa{tornberg2024}, \pa{pangakis2023}).
\item Report per code: precision, recall, F1, $\kappa$ with counts and prevalence, AC1 for rare codes, and the number of validation positives; state that model--human agreement is not human reliability (\pa{harkous2018}, \pa{mcdonald2019}).
\item Ethics: what text is sent to which provider, under which terms, with what retention (\pa{tornberg2024}, \pa{aoir2019}, \pa{mayring2025}).
\end{itemize}
\end{protocolbox}

\begin{thesisbox}[Implications for this thesis (stage L)]
The draft design is consistent with the best-practice sources: the model is a screener with recall-oriented prompts, prompts are staged and frozen on the pilot, a blind validation set is coded before model output is seen, two paraphrased prompts test prompt sensitivity, and absence claims are withheld when recall is low. To close D9 the Decision Log should add: the pre-registered configuration (model, version, prompt, parameters, mapping), acceptance thresholds per code with the rule for rare codes, the number of validation positives below which a code is reviewed manually, and the statement that LLM labels are never used for inductive coding (\pa{mayring2025}). Because the task is deductive coding with a closed codebook and human review of every retained mention, the thesis stays inside the setting where \pa{pangakis2023} and \pa{tornberg2024} report the best evidence; it should still report the multiverse of prompts and models it actually ran.
\end{thesisbox}
